% Ne pas toucher au taux, mieux déduire l'investissement (rapport Sorbon, partie III)
% Compilation depuis fiscalite/is/ : pdflatex note-is.tex (deux passes)
% Matière et statut des sources : notes/is_recherche.md ; chiffres : donnees/is_chiffres.csv
% Calculs : python3 simulation_is.py -> sorties/simulation_is.txt ; figure : python3 is/figures_is.py (depuis fiscalite/)
\documentclass[11pt,a4paper]{article}
% Préambule commun au brief et à la revue de littérature sur la taxe sur la valeur foncière.
% Même composition que le cours (cours/refaire-l-impot.tex) ; figures « nues » produites par
%   FIG_NU=1 python lvt/graphiques_lvt.py
\usepackage[T1]{fontenc}
\usepackage[utf8]{inputenc}
\usepackage{lmodern}
\usepackage[french]{babel}
\usepackage[a4paper,left=2.8cm,right=2.8cm,top=2.6cm,bottom=2.8cm]{geometry}
\usepackage{graphicx}
\usepackage[export]{adjustbox}
\usepackage{booktabs}
\usepackage{tabularx}
\usepackage{array}
\usepackage{xcolor}
\usepackage{titlesec}
\usepackage[labelfont=bf,labelsep=period,font=small,singlelinecheck=false,
            justification=raggedright,skip=6pt]{caption}
\usepackage{fancyhdr}
\usepackage{enumitem}
\usepackage{microtype}
\usepackage{textcomp}
\usepackage{amsmath}
\providecommand{\euro}{\texteuro}
\usepackage[hidelinks]{hyperref}

\definecolor{encre}{HTML}{1A1A19}
\definecolor{bleu}{HTML}{003399}
\definecolor{vert}{HTML}{1F7A4D}
\definecolor{ocre}{HTML}{C79100}
\definecolor{brique}{HTML}{993333}
\definecolor{filet}{HTML}{C8C6BE}

\titleformat{\section}{\normalfont\large\bfseries\color{bleu}}{\thesection.}{0.6em}{}
\titleformat{\subsection}{\normalfont\normalsize\bfseries\color{bleu}}{\thesubsection}{0.6em}{}
\titlespacing*{\section}{0pt}{2.4ex plus 1ex minus .2ex}{1.2ex plus .2ex}
\titlespacing*{\subsection}{0pt}{2ex plus 1ex minus .2ex}{0.8ex plus .2ex}

\pagestyle{fancy}
\fancyhf{}
\renewcommand{\headrulewidth}{0.4pt}
\fancyhead[R]{\small\thepage}

\setlength{\parindent}{0pt}
\setlength{\parskip}{0.62em}
\linespread{1.05}
\renewcommand{\arraystretch}{1.2}
\setlist{itemsep=2pt,topsep=2pt,parsep=0pt,leftmargin=1.4em}

\newenvironment{encadre}[1]
  {\par\medskip\noindent\begin{minipage}{\linewidth}%
   \color{encre}\setlength{\parskip}{0.5em}%
   \hrule height 0.8pt width \linewidth \vspace{6pt}%
   {\footnotesize\bfseries\color{vert}\MakeUppercase{#1}}\par\vspace{-2pt}%
   \small}
  {\par\vspace{4pt}\hrule height 0.4pt width \linewidth\end{minipage}\par\medskip}

\newcommand{\figbloc}[3]{%
  \begin{figure}[htbp]\centering
  \caption{#2}\label{#3}
  \includegraphics[max width=\linewidth,max height=0.40\textheight]{../figures/nu/#1.pdf}
  \IfFileExists{../figures/nu/#1_note.tex}{\par\vspace{4pt}%
    \parbox{\linewidth}{\raggedright\scriptsize\color{encre}\setlength{\parskip}{2pt}%
    \input{../figures/nu/#1_note.tex}}}{}%
  \end{figure}}

\newcommand{\pc}{\,\%}

\renewcommand{\figbloc}[3]{%
  \begin{figure}[!htbp]\centering
  \caption{#2}\label{#3}
  \includegraphics[max width=\linewidth,max height=0.29\textheight]{../figures/nu/#1.pdf}
  \par\vspace{3pt}\parbox{\linewidth}{\raggedright\scriptsize\color{encre}\setlength{\parskip}{2pt}%
    \input{../figures/nu/#1_note.tex}}%
  \end{figure}}
\usepackage{needspace}
\usepackage{etoolbox}
\preto\section{\needspace{6\baselineskip}}
\renewcommand{\topfraction}{0.9}\renewcommand{\bottomfraction}{0.8}
\renewcommand{\textfraction}{0.08}\renewcommand{\floatpagefraction}{0.85}
\fancyhead[L]{\small\itshape Ne pas toucher au taux, mieux déduire l'investissement}

\begin{document}
\renewcommand{\figurename}{Figure}\renewcommand{\tablename}{Tableau}

\thispagestyle{empty}
\vspace*{0.6cm}
{\footnotesize\color{ocre}\bfseries NOTE DE TRAVAIL\quad·\quad NOVEMBRE 2026\par}
\vspace{8pt}
{\fontsize{22}{26}\selectfont\bfseries\color{bleu}Ne pas toucher au taux, mieux déduire
l'investissement\par}
\vspace{6pt}
{\large L'impôt sur les sociétés : déduction immédiate des équipements, financée par la fin du taux réduit
et de trois niches\par}
\vspace{8pt}
{\small Colin Baget. Droit de 2026 ; Tax Foundation, OCDE, Conseil des prélèvements obligatoires ; comptes
nationaux 2025 et statistiques de la DGFiP ; programmes \texttt{fiscalite/is/}.\par}
\vspace{6pt}
{\color{bleu}\hrule height 1.2pt}
\vspace{4pt}

\section*{Introduction}

L'impôt sur les sociétés (IS) n'appelle pas de grande réforme. Ce que lui reprochent les comparaisons
internationales, c'est son taux affiché, le plus élevé de l'OCDE pour les plus grands groupes en 2025 et 2026
avec la contribution exceptionnelle, sa complexité et son instabilité. Ce n'est pas le traitement de
l'investissement, qui est dans la moyenne. Cette note propose de garder le taux normal à 25~\%, de laisser
la contribution exceptionnelle s'éteindre comme prévu, et de faire un seul changement : permettre aux
entreprises de déduire tout de suite la moitié du coût de leurs équipements, puis la totalité quand les
finances publiques le permettent. Le financement vient de l'assiette : la fin du taux réduit de 15~\% des
PME et le resserrement de trois niches dont l'efficacité n'est pas démontrée.

Porter le taux à 28~\% pour payer cette mesure serait une erreur de calendrier autant que de principe. La
déduction immédiate coûte cher les premières années, puis peu : de l'ordre de 2~Md€ par an, moins d'un point
de taux. Une hausse de trois points rapporterait durablement deux à quatre fois plus, en relevant le
paramètre que regardent les groupes qui choisissent où investir et où déclarer leurs bénéfices. La fin du
taux réduit, elle, pèse surtout sur les petites sociétés qui investissent peu, dont beaucoup de sociétés de
professions libérales et de sociétés patrimoniales ; elle coûte au plus 4\,250~€ par an à chacune.

\section{Ce que les comparaisons reprochent à l'IS français}

La Tax Foundation classe la France dernière des 38 pays de l'OCDE en 2025 et en 2026, au total comme pour
l'IS. Ce qui la pénalise est le taux de 36,1~\% des plus grands groupes, contre 24,2~\% en moyenne, quatre taux
différents quand la moyenne est de 2,3, une surtaxe, une boîte à brevets et la troisième subvention implicite
à la recherche la plus généreuse. L'indice est normatif : il pénalise par construction toute incitation
ciblée. Il ne pénalise pas, en revanche, la façon dont la France laisse déduire l'investissement : elle est
15\up{e} sur 38 pour la récupération des coûts, et les déductions d'une machine valent 88~\% de son prix en
valeur actuelle, contre 86~\% en moyenne.

Les taux effectifs de l'OCDE disent la même chose. Le taux moyen effectif, qui compte quand une entreprise
choisit dans quel pays implanter un projet rentable, est de 23,6~\% en France contre 22,2~\% en moyenne
(13\up{e} sur 38). Le taux marginal effectif, qui compte pour décider d'un investissement de plus, est de
15,0~\% contre 16,5~\%, et de 9,8~\% contre 10,4~\% pour les équipements. Ces chiffres de 2025 semblent ne pas
intégrer la contribution exceptionnelle, ce qui n'a pas pu être vérifié. Le Conseil des prélèvements
obligatoires (2023, 2025) et la mission d'information du Sénat (2026) font de l'instabilité le premier
grief : la contribution exceptionnelle, présentée comme temporaire, a été prorogée une fois, et son assiette
est la moyenne de l'IS de deux exercices, dont l'un était clos quand elle a été votée.

La déduction immédiate de l'investissement ne corrige donc pas un retard français. Elle se justifie par ce
qu'elle rapporte en investissement pour chaque euro de recette abandonné.

\section{Pourquoi déduire l'investissement plutôt que baisser le taux}

Une entreprise achète une machine de 100~€ qui dure sept ans. Aujourd'hui, elle en déduit le coût de son
bénéfice par fractions, sur sept ans ; actualisées à 5~\%, ces déductions valent 85 à 90~€ selon que
l'amortissement est linéaire ou dégressif. Si elle peut déduire les 100~€ l'année de l'achat, le total
déduit ne change pas, seulement sa date. Mais ce décalage change la nature de l'impôt. L'État rend 25~€
l'année de l'achat et prélève ensuite 25~\% de ce que la machine rapporte : il devient un associé qui finance
le quart de l'investissement et reçoit le quart des gains. Un investissement qui rapporte tout juste le
rendement normal du capital, celui qu'on obtiendrait en plaçant l'argent, ne paie alors plus d'impôt ; seuls
restent imposés les profits qui dépassent ce rendement (marques, positions de marché, rentes) et les profits
du capital déjà installé.

Une baisse du taux procède à l'inverse : elle allège l'impôt sur tous les profits, y compris ceux des
machines achetées il y a dix ans et ceux des rentes, qui ne dépendent pas de la décision d'investir
aujourd'hui. L'essentiel de son coût ne change donc aucune décision. Chodorow-Reich, Smith, Zidar et Zwick
(2026) le chiffrent sur la réforme américaine de 2017, qui a fait les deux : pour obtenir un point de capital
supplémentaire au bout de dix ans, la baisse du taux coûte chaque année environ 10~\% de la recette de l'IS,
la déduction immédiate 2,5~\%, quatre fois moins.

L'effet sur l'investissement est établi par plusieurs évaluations sur données fiscales. Aux États-Unis, les
déductions immédiates partielles de 2001-2004 et 2008-2010 ont augmenté l'investissement éligible de 10 et
17~\%, deux fois plus dans les petites entreprises, et seulement dans celles qui payaient l'impôt, pour qui la
déduction apporte de la trésorerie tout de suite (Zwick et Mahon, 2017). Au Royaume-Uni, des entreprises
moyennes devenues éligibles à une déduction accélérée ont accru leur taux d'investissement d'environ 11~\%
(Maffini, Xing et Devereux, 2019). Deux réserves comptent. Ces mesures étaient temporaires et ont en partie
avancé des achats ; l'Office for Budget Responsibility estime qu'une déduction temporaire déplace
l'investissement dans le temps sans augmenter durablement le capital, d'où l'importance qu'elle soit
permanente. Et il n'existe aucune évaluation causale du suramortissement français de 2015 ; la seule étude
française, sur celui de 2019-2020, trouve un effet de 12 à 14~\%, difficile à concilier avec ses 2\,390
utilisateurs (DGE, 2026).

L'ordre de grandeur pour la France peut s'écrire avec le coût d'usage du capital, ce que doit rapporter une
machine pour couvrir son financement, son usure et l'impôt :
\[
c = (r+\delta)\,\frac{1-\tau z}{1-\tau},
\]
où $\tau$ est le taux de l'IS et $z$ la valeur actuelle des déductions pour 1~€ investi. Avec $z$ entre 0,85
et 0,90, l'impôt majore le coût d'usage de 3,5 à 5~\% ; avec la déduction immédiate, $z = 1$ et l'impôt ne le
majore plus. Si le stock de capital réagit à peu près proportionnellement au coût d'usage, comme le suggèrent
les estimations britanniques, le stock d'équipements des entreprises imposables augmenterait de quelques
pour cent à long terme. C'est une extrapolation, pas une mesure.

La mesure a trois limites. Combinée à la déduction des intérêts, elle subventionne l'investissement financé
par emprunt, ce que la limitation des charges financières à 30~\% du résultat brut atténue. Elle vaut moins
pour les entreprises en perte, d'autant que l'imputation des déficits est plafonnée. Et elle crée une
frontière entre les actifs éligibles et les autres : on la limite ici aux machines, équipements et
matériels informatiques, ceux dont l'investissement réagit le plus, en laissant de côté les bâtiments, déjà
peu taxés en France (taux marginal effectif de 7,3~\% contre 12,1~\% dans l'OCDE), et les véhicules.

\section{Combien cela coûte, et pourquoi 28~\% n'est pas la réponse}

Les sociétés non financières ont investi 84~Md€ en 2025 en machines, équipements et matériels
informatiques. On suppose que les trois quarts sont le fait de sociétés imposables, qui peuvent utiliser la
déduction tout de suite ; cette part n'est pas mesurée et le coût lui est proportionnel. La première année
(2027), déduire tout l'investissement au lieu de son premier amortissement coûte 14 à 15,5~Md€ (11 à 19~Md€
selon la part imposable). Le coût baisse ensuite chaque année, parce que les machines achetées après la
réforme n'ont plus d'amortissement à déduire : l'État récupère ce qu'il a avancé. Au bout de sept ans, il ne
reste que le coût permanent, 1,4 à 2,1~Md€ par an, qui vient de ce que l'investissement croît et de ce que
l'avance a une valeur. Mesuré en valeur actuelle, le coût d'une année d'investissement est de 1,8 à
2,6~Md€, soit 0,7 à 1 point de taux (figure~\ref{fig:cout}). Ces chiffres ne comptent pas l'investissement
supplémentaire suscité, qui augmente le coût la première année et la recette ensuite.

\figbloc{is_01_cout_deduction}{Coût annuel de la déduction immédiate des équipements et ressources pour la
financer}{fig:cout}

Une hausse du taux normal de 25 à 28~\% rapporterait 7,5~Md€ à bénéfices inchangés, et 4,4~Md€ si les
bénéfices déclarés baissent de 1,5~\% par point de taux, l'ordre de grandeur que trouve la littérature sur
les filiales de multinationales (Beer, de Mooij et Liu, 2020). C'est trop la plupart des années : elle
financerait le coût initial, concentré sur cinq ans, par une hausse permanente, deux à quatre fois supérieure
au coût de long terme. Elle relèverait surtout le paramètre qui compte pour les décisions les plus mobiles :
où implanter un projet très rentable, qui dépend du taux moyen effectif, et où déclarer les bénéfices, qui
dépend du taux affiché. L'avantage de la déduction immédiate est de ne pas toucher à l'imposition des rentes ;
le financer par le taux reviendrait à alourdir précisément les profits les plus mobiles.

Le financement par l'assiette suffit pour le régime permanent (tableau~\ref{tab:assiette}). Les quatre
mesures retenues rapportent 3,5 à 4,4~Md€ par an, soit 1,4 à 1,7 point de taux. La première est la fin du
taux réduit, traitée au §~4. Le crédit d'impôt recherche (8~Md€) produit environ un euro de recherche par
euro de crédit, avec des effets concentrés sur les PME : on le garde, mais on l'apprécie au niveau du groupe
plutôt que de chaque filiale, comme le propose le CPO (2025).
Les exonérations zonées sont jugées « peu efficientes » par le Conseil des prélèvements obligatoires. Le
coût de la boîte à brevets a doublé depuis 2021, et son efficacité est « mise en doute » par le même
Conseil.

\begin{table}[!htbp]\centering\small
\caption{Ressources tirées de l'assiette}\label{tab:assiette}
\begin{tabularx}{\linewidth}{@{}>{\raggedright\arraybackslash}X>{\raggedleft\arraybackslash}p{2.4cm}>{\raggedright\arraybackslash}p{5.2cm}@{}}
\toprule
Mesure & Md€ par an & Source du chiffrage \\
\midrule
Fin du taux réduit de 15~\% & 2,2 à 2,7 & Chiffre officiel de 2019 ; calcul statique sur 2022 \\
Crédit d'impôt recherche apprécié au niveau du groupe & 0,6 & CPO (2025) \\
Fin des exonérations zonées, hors outre-mer & 0,5 & CPO (2025) \\
Avantage de la boîte à brevets réduit de moitié (taux de 17,5~\%) & 0,2 à 0,6 & Calcul ; dépend de la fuite des brevets \\
\midrule
Total & 3,5 à 4,4 & Soit 1,4 à 1,7 point de taux normal \\
\bottomrule
\end{tabularx}
\end{table}

Ce qui ne tient pas d'emblée, c'est la transition. Une déduction immédiate totale dès 2027 creuserait le
solde cumulé de 18 à 40~Md€ au bout de quatre à six ans, regagnés en dix à vingt-deux ans : l'État prête aux
entreprises, et le prêt se rembourse, mais la France n'a pas aujourd'hui la marge budgétaire pour le
consentir. On commence donc par une déduction immédiate de 50~\%, le reste étant amorti comme aujourd'hui, et
permanente. Elle coûte 7 à 7,8~Md€ la première année ; avec l'élargissement de l'assiette, le solde cumulé
descend au plus de 3 à 10,5~Md€ et redevient positif en quatre à neuf ans. Elle réduit de moitié l'écart entre
la valeur actuelle des déductions et le prix de la machine. Le passage à 100~\% vient dans un second temps et
a sa propre bosse de coût, plus petite : la montée par paliers étale la dépense, elle ne la supprime pas.

\section{La fin du taux réduit : qui y perd}

Les sociétés dont le chiffre d'affaires est inférieur à 10~M€ et dont le capital est détenu à 75~\% par des
personnes physiques paient 15~\% au lieu de 25~\% sur les 42\,500 premiers euros de bénéfice (art.~219 du
CGI). Le Conseil d'État a jugé en 2025 que le chiffre d'affaires s'apprécie au niveau du groupe, ce qui ferme
l'usage par les filiales de grands groupes. Le dispositif a coûté 2,2~Md€ en 2019 pour 760\,000 sociétés ; le
bénéfice imposé à 15~\% en 2022 (27~Md€) donne environ 2,7~Md€.

Le taux réduit ne finance ni l'investissement ni la croissance des entreprises. Pour une société dont le
bénéfice dépasse le plafond, il revient à verser chaque année 4\,250~€, quoi qu'elle fasse : la somme ne
dépend d'aucune de ses décisions, sauf celle de rester petite ou de se découper en plusieurs sociétés. Le
CPO en demande la suppression depuis 2016 : les PME ne sont pas moins rentables que les grandes entreprises,
on observe un pic de sociétés juste sous le plafond de bénéfice, signe que des résultats sont pilotés pour
rester en dessous, et aucun effet sur l'émergence d'entreprises de taille intermédiaire n'a été trouvé. La
seule étude qui trouve un effet réel, une hausse de 4~\% des actifs corporels lors de son introduction en
2001-2003, est fragile et non publiée (Bernini et Treibich, 2014). Le meilleur argument en sa faveur, aider
les petites entreprises à se financer sur leurs bénéfices quand les banques leur prêtent mal, est mieux servi
par la déduction immédiate, qui donne de la trésorerie aux seules entreprises qui investissent.

Qui y perd ? Toutes les sociétés éligibles, au plus 4\,250~€ par an chacune, soit au plus 10~\% de leur
bénéfice ; celles dont le bénéfice est inférieur au plafond voient leur IS augmenter des deux tiers. La perte
est sèche pour celles qui investissent peu en équipements : sociétés de professions libérales, sociétés de
conseil et de services, petites holdings et sociétés civiles immobilières soumises à l'IS. Aucune
statistique ne ventile la base imposée à 15~\% par type de société ; on peut seulement noter que, sur
2,4~millions de redevables de l'IS, 490\,000 relèvent des activités immobilières et 190\,000 des activités
financières, où sont classées les holdings, et supposer, sans pouvoir le mesurer, qu'une part notable du
taux réduit leur profite. Une petite société industrielle qui investit est en partie compensée : la
déduction immédiate lui avance l'impôt sur ses achats, 4\,600~€ la première année pour 20\,000~€
d'équipements, mais son gain durable ne dépasse la perte que si elle investit plus de 110\,000~€ par an avec une déduction
totale, le double avec la déduction de moitié.
L'opération transfère donc de l'IS des petites sociétés qui investissent peu vers les entreprises de toutes
tailles qui investissent. Pour l'associé qui se fait distribuer le bénéfice, le taux cumulé de l'IS et du PFU
sur les 42\,500 premiers euros passe de 41,7 à 48,6~\%, le taux de toutes les autres sociétés, ce qui supprime
une raison purement fiscale de passer à l'IS.

\section{Ce qu'on ne fait pas}

On n'introduit pas de barème progressif de l'IS. La taille d'une entreprise ne dit rien de la richesse de
ses propriétaires : une société de chirurgiens appartient à des ménages aisés, une grande société cotée en
partie à des fonds de retraite et à des épargnants modestes. La progressivité se règle dans l'impôt des
personnes, au barème fusionné et au PFU. Un barème par tranches de bénéfice multiplierait en outre les pics
sous les seuils et les découpages en sociétés que le taux réduit produit déjà.

On n'introduit pas de taux minimal national. Le minimum mondial de 15~\% s'applique déjà aux groupes dont
le chiffre d'affaires consolidé dépasse 750~M€ (loi de finances pour 2024, art.~33 ; 0,5~Md€ attendus en
2026). Un minimum assis sur le taux effectif frapperait précisément les entreprises qui investissent
beaucoup une année donnée, puisque la déduction immédiate abaisse leur taux effectif cette année-là : il
annulerait la mesure proposée.

On laisse la contribution exceptionnelle s'éteindre, comme le demandent le CPO et le Sénat : une surtaxe
dont l'assiette est en partie rétroactive et qui est renouvelée d'année en année fait payer les
investissements déjà faits et rend incertain le rendement de ceux qu'on envisage. Les recettes de référence
de cette note sont celles du droit permanent, à 25~\%. Enfin, la contribution sociale de 3,3~\% pourrait être
fondue dans le taux normal, ce qui supprimerait un taux et un seuil sans changer la recette ; c'est une
simplification secondaire.

L'ensemble ne change pas la charge des grandes fortunes de façon notable : les grands groupes paient un peu
plus au titre du crédit d'impôt recherche et des brevets, et moins sur leurs investissements. Il n'a pas
d'effet propre sur les familles.

\vspace{6pt}
{\scriptsize\textbf{Références.} Beer, de Mooij et Liu (2020), « International Corporate Tax Avoidance: A
Review of the Channels, Magnitudes, and Blind Spots », \emph{Journal of Economic Surveys} 34(3). Bernini et
Treibich (2014), document de travail FREIT n°~759. Chodorow-Reich, Smith, Zidar et Zwick (2026), « Tax
Policy and Investment in a Global Economy », NBER WP 32180 (version de janvier 2026). Conseil des
prélèvements obligatoires (2016, 2023, 2025), rapports sur l'IS, sur les PME et grandes entreprises, et sur
le cadre fiscal de l'industrie. Conseil d'État, 13 mars 2025, n°~481538. DGE (2026), « Le suramortissement :
un levier efficace pour la modernisation industrielle des PME françaises ? », document de travail n°~1.
DGFiP, \emph{Statistiques} n°~21 (2024). IFS (2023), \emph{Green Budget}, ch.~10. Insee, comptes nationaux
annuels, base 2020 (2026). Maffini, Xing et Devereux (2019), « The Impact of Investment Incentives: Evidence
from UK Corporation Tax Returns », \emph{AEJ: Economic Policy} 11(3). OCDE, \emph{Corporate Tax Statistics}
(2025). Sénat (2026), mission d'information sur le poids des prélèvements obligatoires sur les entreprises.
Tax Foundation, \emph{International Tax Competitiveness Index} 2025 et 2026. Zwick et Mahon (2017), « Tax
Policy and Heterogeneous Investment Behavior », \emph{American Economic Review} 107(1). Détail et statut de
lecture : \texttt{is/notes/is\_recherche.md}.\par}

\end{document}
